\chapter{Diagrama de Casos de Uso}
\label{cap:casosuso}

\section{Diagrama gráfico de alto nivel}
\label{sec:cu_altonivel}

\figuradiagrama{1.0}{01-casos-uso-alto-nivel.png}
{Diagrama de casos de uso de alto nivel del sistema BeeTracer}
{Diagrama de casos de uso de alto nivel del sistema BeeTracer, organizado por módulos.
Fuente PlantUML: \texttt{diagramas/fuente/01-casos-uso-alto-nivel.puml}.}
{fig:cu_altonivel}

\begin{nota}
El caso de uso \textbf{CU-22 Autenticar usuario} es incluido (\ster{include}) por la totalidad de
los casos de uso del sistema. Se omite del diagrama general por legibilidad.
\end{nota}

\section{Diagrama gráfico expandido: subsistema de ejecución en terreno}
\label{sec:cu_terreno}

El módulo \textbf{M4. Ejecutar operaciones} concentra el valor del sistema y la mayor densidad de
relaciones \ster{include} y \ster{extend}, por lo que se presenta en un diagrama aparte con mayor
nivel de detalle.

\figuradiagrama{1.0}{02-casos-uso-ejecucion-terreno.png}
{Casos de uso del subsistema de ejecución en terreno}
{Casos de uso del subsistema de ejecución en terreno (aplicación móvil).
Fuente PlantUML: \texttt{diagramas/fuente/02-casos-uso-ejecucion-terreno.puml}.}
{fig:cu_terreno}

\section{Casos de uso de alto nivel (narrativos)}
\label{sec:cu_narrativos}

La Tabla~\ref{tab:cu_altonivel} presenta los casos de uso en formato breve según Larman:
identificador, actores, tipo y descripción en un párrafo.

{\estilotabla
\begin{longtable}{L{0.06\textwidth} L{0.18\textwidth} L{0.14\textwidth} L{0.11\textwidth} L{0.41\textwidth}}

\caption{Casos de uso de alto nivel del sistema BeeTracer}
\label{tab:cu_altonivel}
\\

\toprule
\textbf{ID} & \textbf{Caso de uso} & \textbf{Actores} & \textbf{Tipo} & \textbf{Descripción} \\
\midrule
\endfirsthead

\toprule
\textbf{ID} & \textbf{Caso de uso} & \textbf{Actores} & \textbf{Tipo} & \textbf{Descripción} \\
\midrule
\endhead

CU-01 & Gestionar usuarios, roles y permisos & Administrador del depósito & Primario, esencial &
El Administrador crea, modifica o desactiva cuentas de usuario, define roles y les asigna permisos
granulares sobre las funcionalidades del sistema. El sistema valida la unicidad de las credenciales
y persiste la configuración. \\
\midrule

CU-02 & Gestionar clientes & Administrador, Ejecutivo comercial & Primario, esencial &
Se registra o actualiza la ficha de una empresa importadora, incluyendo su identificación
tributaria y los correos destinatarios a los que se enviarán sus reportes. El sistema valida los
datos y los deja disponibles para asociarlos a las líneas de carga. \\
\midrule

CU-03 & Gestionar trabajadores y cuadrillas & Administrador del depósito & Primario, esencial &
Se registran los trabajadores del depósito y se los agrupa en cuadrillas de cuatro a cinco
personas, designando quiénes pueden actuar como jefe tarjador. El sistema deja las cuadrillas
disponibles para su asignación a faenas. \\
\midrule

CU-04 & Gestionar tipos de bulto y parámetros & Administrador del depósito & Primario, secundario &
Se parametrizan los tipos de bulto y unidades de medida con los que se registrará la carga,
ajustando el sistema a la operación particular del depósito. \\
\midrule

CU-05 & Importar manifiesto electrónico (XML) & Jefe de faena, Servicio Nacional de Aduanas &
Primario, esencial &
El Jefe de faena solicita la importación del manifiesto de una nave. El sistema se conecta al
\este{web service} de Aduanas, descarga el XML, valida su esquema, lo interpreta y persiste
contenedores, líneas de carga, consignatarios y cantidades declaradas, quedando disponible el
contenido del contenedor antes de su arribo. \\
\midrule

CU-06 & Recibir carga vía integración API & Sistema del cliente (API) & Primario, esencial &
El sistema informático del cliente se conecta a la API de BeeTracer y transfiere directamente la
información de la carga, que el sistema valida y persiste. \\
\midrule

CU-07 & Registrar carga en modalidad tarja ciega & Jefe tarjador & Primario, esencial &
Cuando no existe información previa de la carga, la aplicación móvil arranca en blanco y el
operario ingresa manualmente la descripción de cada lote leyendo las viñetas y facturas adheridas a
las cajas a medida que las descarga. \\
\midrule

CU-08 & Programar faena de desconsolidado & Jefe de faena & Primario, esencial &
El Jefe de faena consulta los contenedores por arribar, selecciona uno y define fecha y franja
horaria para su desconsolidado. El sistema valida la disponibilidad y registra la faena en estado
programada. \\
\midrule

CU-09 & Asignar cuadrilla y jefe tarjador & Jefe de faena & Primario, esencial &
Sobre una faena programada, el Jefe de faena selecciona la cuadrilla que la ejecutará y designa al
integrante responsable de operar la aplicación móvil. El sistema valida disponibilidad y publica la
faena en el dispositivo del operario designado. \\
\midrule

CU-10 & Consultar programación (móvil) & Jefe tarjador, Jefe de faena & Primario, esencial &
Al iniciar sesión en la aplicación móvil, el operario visualiza las faenas que le fueron asignadas,
con su contenedor, horario y contenido declarado. \\
\midrule

CU-11 & Registrar apertura de contenedor & Jefe tarjador & Primario, esencial &
Antes de abrir, el operario fotografía puertas, candado y sello, y digita el número de sello. El
sistema contrasta el número contra el declarado, registra el estado del sello y solo entonces
habilita la etapa siguiente. \\
\midrule

CU-12 & Registrar desconsolidado y separación por cliente & Jefe tarjador, Cuadrilla &
Primario, esencial &
A medida que la cuadrilla extrae la mercancía, el operario la agrupa por cliente, registra tipo de
bulto y cantidad recibida de cada lote y lo fotografía. El sistema contrasta la cantidad recibida
contra la declarada y registra las diferencias. \\
\midrule

CU-13 & Registrar incidencia & Jefe tarjador & Primario, esencial &
Ante una caja dañada o una diferencia de cantidad, el operario separa la unidad afectada, clasifica
la incidencia y le toma fotografías exclusivas. El sistema asocia la incidencia al lote y al
cliente correspondiente. \\
\midrule

CU-14 & Registrar cierre de faena & Jefe tarjador & Primario, esencial &
El operario fotografía el contenedor completamente vacío. El sistema verifica que toda la evidencia
obligatoria esté completa y sincronizada, cierra la faena y dispara la generación del reporte. \\
\midrule

CU-15 & Capturar y subir evidencia fotográfica & Jefe tarjador & Secundario, esencial &
El sistema captura la imagen desde la cámara, la comprime, la sella con fecha, hora y
geoposición, y la sube individualmente al repositorio inmediatamente después de tomada. \\
\midrule

CU-16 & Sincronizar datos (modo offline) & --- \este{(disparado por el sistema)} &
Secundario, esencial &
Sin conectividad, el sistema persiste las capturas en la base de datos local del dispositivo y las
encola; al detectar señal, sincroniza la cola de pendientes y confirma la subida. \\
\midrule

CU-17 & Generar reporte PDF de la faena & --- \este{(disparado por el sistema)} &
Secundario, esencial &
Al cerrarse una faena, el sistema compila automáticamente datos y fotografías en un PDF con
encabezado, desglose por cliente e incidencias destacadas en rojo, y lo deja en estado borrador. \\
\midrule

CU-18 & Revisar y aprobar reporte & Supervisor de reportes & Primario, esencial &
El Supervisor abre el borrador en la plataforma web, lo revisa visualmente y lo aprueba. El sistema
registra quién aprobó y cuándo, y habilita su publicación. \\
\midrule

CU-19 & Solicitar recaptura de fotografía & Supervisor de reportes & Primario, esencial &
Ante una fotografía defectuosa, el Supervisor la rechaza indicando el motivo. El sistema la marca
como rechazada, notifica al jefe tarjador y, tras recibir el reemplazo, regenera y versiona el
reporte. \\
\midrule

CU-20 & Enviar reporte al cliente & Supervisor de reportes, Servidor de correo &
Primario, esencial &
Aprobado el reporte, el sistema obtiene los destinatarios del cliente y despacha el PDF por correo
electrónico, registrando el acuse de envío y reintentando ante fallo. \\
\midrule

CU-21 & Consultar faenas históricas y trazabilidad &
Jefe de faena, Supervisor, Cliente del depósito & Primario, esencial &
El usuario consulta faenas en proceso o históricas, filtrando por evento, cliente o fecha, y
recupera sus imágenes y su reporte PDF. \\
\midrule

CU-22 & Autenticar usuario & Todos los actores humanos & Secundario, esencial &
El sistema valida las credenciales del usuario, establece su sesión y determina los permisos
asociados a su rol. \\

\bottomrule

\end{longtable}
}

\section{Casos de uso expandidos}
\label{sec:cu_expandidos}

\subsection{CU-05 --- Importar manifiesto electrónico (XML)}
\label{subsec:cu05}

{\estilotabla
\begin{longtable}{L{0.25\textwidth} L{0.69\textwidth}}

\caption{CU-05 Importar manifiesto electrónico (XML): descripción expandida}
\label{tab:cu05_desc}
\\

\toprule
\textbf{Campo} & \textbf{Contenido} \\
\midrule
\endfirsthead

\toprule
\textbf{Campo} & \textbf{Contenido} \\
\midrule
\endhead

\textbf{Caso de uso} & CU-05 Importar manifiesto electrónico (XML) \\
\midrule
\textbf{Actores} & Jefe de faena \este{(iniciador)}, Servicio Nacional de Aduanas
\este{(sistema externo)} \\
\midrule
\textbf{Propósito} & Conocer el detalle exacto de la carga que transporta un contenedor
\textbf{antes de que la nave recale en el puerto}, precargando el sistema con información
oficial. \\
\midrule
\textbf{Tipo} & Primario $\cdot$ Esencial \\
\midrule
\textbf{Referencias cruzadas} & Funciones R2.1, R2.2, R2.3, R2.4, R2.7 $\cdot$
Casos de uso CU-06, CU-07, CU-08, CU-22 \\
\midrule
\textbf{Precondiciones} & El Jefe de faena está autenticado y posee el permiso de ingesta de carga.
El sistema tiene configurada y vigente la conexión al \este{web service} de Aduanas. \\
\midrule
\textbf{Postcondiciones} & El manifiesto, sus contenedores y sus líneas de carga quedan persistidos
y asociados a los clientes correspondientes, disponibles para programar faenas. \\

\bottomrule

\end{longtable}
}

{\estilotabla
\begin{longtable}{C{0.04\textwidth} L{0.44\textwidth} L{0.44\textwidth}}

\caption{CU-05: curso normal de los eventos}
\label{tab:cu05_normal}
\\

\toprule
\textbf{\#} & \textbf{Acción de los actores} & \textbf{Respuesta del sistema} \\
\midrule
\endfirsthead

\toprule
\textbf{\#} & \textbf{Acción de los actores} & \textbf{Respuesta del sistema} \\
\midrule
\endhead

1 & El Jefe de faena ingresa a la opción de importación de manifiestos e indica nave, viaje y
fecha. & \\
\midrule
2 & & El sistema se conecta al \este{web service} del Servicio Nacional de Aduanas. \\
\midrule
3 & & El sistema solicita y descarga el documento XML correspondiente. \\
\midrule
4 & & El sistema valida el esquema del XML recibido. \\
\midrule
5 & & El sistema interpreta el documento y crea el registro del manifiesto (número, nave, viaje,
fecha de recepción). \\
\midrule
6 & & Por cada contenedor declarado, el sistema crea el registro del contenedor con su sigla,
número, tipo, tamaño y número de sello declarado. \\
\midrule
7 & & Por cada línea de carga de cada contenedor, el sistema registra el conocimiento de embarque,
la descripción, las marcas, la cantidad declarada, el peso y el tipo de bulto, asociándola al
cliente consignatario. \\
\midrule
8 & & El sistema persiste la información y marca el origen de los datos como
\estado{MANIFIESTO\_XML}. \\
\midrule
9 & & El sistema presenta un resumen de la importación: cantidad de contenedores y de líneas de
carga incorporadas. \\
\midrule
10 & El Jefe de faena verifica el resumen y continúa con la programación de faenas (CU-08). & \\

\bottomrule

\end{longtable}
}

{\estilotabla
\begin{longtable}{C{0.06\textwidth} L{0.38\textwidth} L{0.48\textwidth}}

\caption{CU-05: cursos alternos}
\label{tab:cu05_alt}
\\

\toprule
\textbf{Paso} & \textbf{Situación} & \textbf{Tratamiento} \\
\midrule
\endfirsthead

\toprule
\textbf{Paso} & \textbf{Situación} & \textbf{Tratamiento} \\
\midrule
\endhead

2 & El \este{web service} de Aduanas no responde o la conexión falla. &
El sistema informa la indisponibilidad del servicio y ofrece reintentar. El Jefe de faena puede
optar por CU-06 (integración API) o declarar tarja ciega (CU-07). \\
\midrule

3 & No existe manifiesto para la nave y viaje indicados. &
El sistema informa que el manifiesto aún no ha sido transmitido por la nave y sugiere reintentar
más cerca de la fecha de arribo. \\
\midrule

4 & El XML no cumple el esquema esperado. &
El sistema rechaza el documento, registra el error en la bitácora y notifica al Jefe de faena. No
se persiste información parcial. \\
\midrule

7 & Una línea de carga referencia un consignatario que no existe como cliente registrado. &
El sistema crea el cliente en estado \este{pendiente de completar} y advierte al Jefe de faena para
que complete su ficha (CU-02) antes de emitir reportes. \\
\midrule

5--8 & El manifiesto ya fue importado previamente. &
El sistema detecta el número de manifiesto duplicado y ofrece \textbf{actualizar} la información
existente en lugar de duplicarla. \\

\bottomrule

\end{longtable}
}

\subsection{CU-08 --- Programar faena de desconsolidado}
\label{subsec:cu08}

{\estilotabla
\begin{longtable}{L{0.25\textwidth} L{0.69\textwidth}}

\caption{CU-08 Programar faena de desconsolidado: descripción expandida}
\label{tab:cu08_desc}
\\

\toprule
\textbf{Campo} & \textbf{Contenido} \\
\midrule
\endfirsthead

\toprule
\textbf{Campo} & \textbf{Contenido} \\
\midrule
\endhead

\textbf{Caso de uso} & CU-08 Programar faena de desconsolidado \\
\midrule
\textbf{Actores} & Jefe de faena \este{(iniciador)} \\
\midrule
\textbf{Propósito} & Planificar la agenda de apertura de contenedores, asignando franja horaria y
equipo de trabajo, de modo que la operación del patio se ejecute de forma ordenada y sin tiempos
muertos. \\
\midrule
\textbf{Tipo} & Primario $\cdot$ Esencial \\
\midrule
\textbf{Referencias cruzadas} & Funciones R3.1 a R3.6 $\cdot$ Casos de uso CU-05, CU-09, CU-10,
CU-22 \\
\midrule
\textbf{Precondiciones} & Existe al menos un contenedor con información de carga cargada y sin
faena asignada. Existen cuadrillas conformadas y trabajadores habilitados. \\
\midrule
\textbf{Postcondiciones} & La faena queda registrada en estado \estado{PROGRAMADA}, con contenedor,
horario, cuadrilla y jefe tarjador asignados, y visible en la aplicación móvil del operario
designado. \\

\bottomrule

\end{longtable}
}

{\estilotabla
\begin{longtable}{C{0.04\textwidth} L{0.44\textwidth} L{0.44\textwidth}}

\caption{CU-08: curso normal de los eventos}
\label{tab:cu08_normal}
\\

\toprule
\textbf{\#} & \textbf{Acción de los actores} & \textbf{Respuesta del sistema} \\
\midrule
\endfirsthead

\toprule
\textbf{\#} & \textbf{Acción de los actores} & \textbf{Respuesta del sistema} \\
\midrule
\endhead

1 & El Jefe de faena abre el módulo de programación. & \\
\midrule
2 & & El sistema lista los contenedores por arribar que aún no tienen faena asignada, con su
contenido declarado y fecha estimada de arribo. \\
\midrule
3 & El Jefe de faena selecciona un contenedor. & \\
\midrule
4 & & El sistema muestra el detalle del contenedor y sus líneas de carga por cliente. \\
\midrule
5 & El Jefe de faena define la fecha y la franja horaria de la faena (p. ej., 10:00 a 12:00). & \\
\midrule
6 & & El sistema valida que no exista conflicto de agenda para esa franja en el patio. \\
\midrule
7 & & El sistema crea la faena, le asigna folio y la deja en estado \estado{PROGRAMADA}. \\
\midrule
8 & & El sistema lista las cuadrillas disponibles en esa fecha y horario. \\
\midrule
9 & El Jefe de faena selecciona una cuadrilla y designa al jefe tarjador responsable
\este{(CU-09)}. & \\
\midrule
10 & & El sistema valida que el trabajador designado esté habilitado como jefe tarjador y
disponible. \\
\midrule
11 & & El sistema asocia cuadrilla y jefe tarjador a la faena y persiste los cambios. \\
\midrule
12 & & El sistema publica la faena en la aplicación móvil del operario designado y confirma la
programación. \\

\bottomrule

\end{longtable}
}

{\estilotabla
\begin{longtable}{C{0.06\textwidth} L{0.38\textwidth} L{0.48\textwidth}}

\caption{CU-08: cursos alternos}
\label{tab:cu08_alt}
\\

\toprule
\textbf{Paso} & \textbf{Situación} & \textbf{Tratamiento} \\
\midrule
\endfirsthead

\toprule
\textbf{Paso} & \textbf{Situación} & \textbf{Tratamiento} \\
\midrule
\endhead

2 & No hay contenedores pendientes. &
El sistema informa que no existen contenedores por programar y ofrece importar un manifiesto
(CU-05). \\
\midrule

6 & La franja horaria está ocupada o excede la capacidad del patio. &
El sistema informa el conflicto, muestra las faenas que colisionan y solicita una nueva franja. \\
\midrule

8 & No hay cuadrillas disponibles en el horario indicado. &
El sistema advierte la falta de disponibilidad y permite guardar la faena sin cuadrilla, en estado
\estado{PROGRAMADA} pendiente de asignación. \\
\midrule

10 & El trabajador designado no está habilitado o ya está asignado a otra faena en el mismo
horario. &
El sistema rechaza la designación e indica el motivo, solicitando elegir otro integrante. \\
\midrule

12 & El operario no tiene conectividad en ese momento. &
La faena queda encolada y se descarga en el dispositivo la próxima vez que la aplicación
sincronice. \\
\midrule

--- & El contenedor cambia de fecha de arribo o se anula la operación. &
El Jefe de faena reprograma o anula la faena (R3.6); el sistema deja registro del cambio y notifica
al operario asignado. \\

\bottomrule

\end{longtable}
}

\subsection{CU-12 --- Registrar desconsolidado y separación por cliente}
\label{subsec:cu12}

\begin{nota}
\textbf{Caso de uso central del sistema.} Concentra la propuesta de valor de BeeTracer: la
generación de evidencia verificable en el momento y lugar exactos en que la mercancía cambia de
custodia.
\end{nota}

{\estilotabla
\begin{longtable}{L{0.25\textwidth} L{0.69\textwidth}}

\caption{CU-12 Registrar desconsolidado y separación por cliente: descripción expandida}
\label{tab:cu12_desc}
\\

\toprule
\textbf{Campo} & \textbf{Contenido} \\
\midrule
\endfirsthead

\toprule
\textbf{Campo} & \textbf{Contenido} \\
\midrule
\endhead

\textbf{Caso de uso} & CU-12 Registrar desconsolidado y separación por cliente \\
\midrule
\textbf{Actores} & Jefe tarjador \este{(iniciador)}, Cuadrilla, Repositorio de fotografías
\este{(sistema externo)} \\
\midrule
\textbf{Propósito} & Documentar fotográfica y cuantitativamente la extracción de la mercancía del
contenedor y su separación por cliente destinatario, dejando constancia irrefutable de qué salió,
en qué cantidad y en qué estado. \\
\midrule
\textbf{Tipo} & Primario $\cdot$ Esencial \\
\midrule
\textbf{Referencias cruzadas} & Funciones R4.6 a R4.11, R5.1 a R5.7 $\cdot$ Casos de uso CU-07,
CU-10, CU-11, CU-13, CU-14, CU-15, CU-16, CU-17 \\
\midrule
\textbf{Precondiciones} & La faena está en estado \estado{EN\_EJECUCION}. Se completó CU-11
(apertura registrada y sello verificado). El operario tiene sesión activa en la aplicación
móvil. \\
\midrule
\textbf{Postcondiciones} & Cada lote extraído queda registrado con su cliente, tipo de bulto,
cantidad recibida y fotografía asociada. Las diferencias respecto de lo declarado quedan calculadas
y, cuando corresponde, escaladas a incidencia. Toda la evidencia queda subida o encolada para
sincronización. \\

\bottomrule

\end{longtable}
}

{\estilotabla
\begin{longtable}{C{0.04\textwidth} L{0.44\textwidth} L{0.44\textwidth}}

\caption{CU-12: curso normal de los eventos}
\label{tab:cu12_normal}
\\

\toprule
\textbf{\#} & \textbf{Acción de los actores} & \textbf{Respuesta del sistema} \\
\midrule
\endfirsthead

\toprule
\textbf{\#} & \textbf{Acción de los actores} & \textbf{Respuesta del sistema} \\
\midrule
\endhead

1 & El Jefe tarjador selecciona la faena en curso e ingresa a la etapa de separación. & \\
\midrule
2 & & El sistema abre la etapa \estado{SEPARACION}, registra su hora de inicio y presenta las
líneas de carga declaradas agrupadas por cliente. \\
\midrule
3 & La cuadrilla extrae mercancía del contenedor y la agrupa en el piso por cliente
destinatario. & \\
\midrule
4 & El Jefe tarjador selecciona el cliente, el tipo de bulto e ingresa la cantidad recibida del
lote. & \\
\midrule
5 & & El sistema contrasta la cantidad recibida contra la cantidad declarada en el manifiesto y
calcula la diferencia. \\
\midrule
6 & & El sistema solicita la captura fotográfica obligatoria del lote. \\
\midrule
7 & El Jefe tarjador toma la fotografía del lote. & \\
\midrule
8 & & El sistema comprime la imagen, la sella con fecha, hora y geoposición, y la asocia a la etapa
y al lote \este{(CU-15)}. \\
\midrule
9 & & El sistema encola la fotografía y la sube individualmente al repositorio en la nube. \\
\midrule
10 & & El sistema confirma el registro del lote y actualiza el avance de la faena. \\
\midrule
11 & Se repiten los pasos 3 a 10 por cada lote extraído del contenedor. & \\
\midrule
12 & El Jefe tarjador indica que el contenedor quedó vacío. & \\
\midrule
13 & & El sistema verifica que todos los clientes declarados tengan al menos un lote registrado y
su evidencia asociada, cierra la etapa \estado{SEPARACION} y habilita la etapa de cierre
\este{(CU-14)}. \\

\bottomrule

\end{longtable}
}

{\estilotabla
\begin{longtable}{C{0.06\textwidth} L{0.38\textwidth} L{0.48\textwidth}}

\caption{CU-12: cursos alternos}
\label{tab:cu12_alt}
\\

\toprule
\textbf{Paso} & \textbf{Situación} & \textbf{Tratamiento} \\
\midrule
\endfirsthead

\toprule
\textbf{Paso} & \textbf{Situación} & \textbf{Tratamiento} \\
\midrule
\endhead

2 & No existe información declarada de la carga (\textbf{tarja ciega}, CU-07). &
La aplicación arranca en blanco y el operario debe ingresar manualmente la descripción, el cliente
y la cantidad de cada lote leyendo viñetas y facturas adheridas a las cajas. El sistema omite el
contraste del paso 5 y marca el origen de datos como \estado{TARJA\_CIEGA}. \\
\midrule

5 & La cantidad recibida difiere de la declarada (faltante o sobrante). &
El sistema levanta automáticamente una incidencia de tipo \estado{FALTANTE} o \estado{SOBRANTE} y
exige fotografías exclusivas del hecho \este{(CU-13)}. \\
\midrule

5 & El operario detecta mercancía dañada o mal embalada. &
El operario separa la unidad afectada y registra una incidencia de tipo \estado{DANO} o
\estado{MAL\_EMBALADO}, con fotografías exclusivas \este{(CU-13)}. La mercancía quedará destacada
en rojo en el reporte. \\
\midrule

7 & La cámara del dispositivo falla o la fotografía sale inutilizable. &
El sistema permite descartar y repetir la captura. \textbf{No permite avanzar sin evidencia.} \\
\midrule

9 & No hay conectividad en el patio. &
El sistema persiste la fotografía en la base de datos local del dispositivo y la mantiene en la
cola de pendientes; el operario continúa trabajando con normalidad. Al recuperar señal, el
sincronizador sube automáticamente los pendientes \este{(CU-16)}. \\
\midrule

9 & La subida de una fotografía falla por caída de la conexión. &
El sistema marca la fotografía como pendiente y reintenta la subida, sin bloquear la operación. \\
\midrule

13 & Un cliente declarado no tiene lote registrado. &
El sistema advierte la omisión y solicita confirmación explícita del operario antes de permitir el
cierre, generando una incidencia de tipo \estado{FALTANTE} por la totalidad del lote. \\
\midrule

--- & Se detecta mercancía no declarada o presuntamente ilícita. &
El operario registra una incidencia de tipo \estado{SOBRANTE} con fotografías, y el procedimiento
continúa según los protocolos de fiscalización del terminal (Aduana, PDI), fuera del alcance del
sistema. \\

\bottomrule

\end{longtable}
}

\subsection{CU-13 --- Registrar incidencia}
\label{subsec:cu13}

{\estilotabla
\begin{longtable}{L{0.25\textwidth} L{0.69\textwidth}}

\caption{CU-13 Registrar incidencia: descripción expandida}
\label{tab:cu13_desc}
\\

\toprule
\textbf{Campo} & \textbf{Contenido} \\
\midrule
\endfirsthead

\toprule
\textbf{Campo} & \textbf{Contenido} \\
\midrule
\endhead

\textbf{Caso de uso} & CU-13 Registrar incidencia (daño, faltante, sobrante o mal embalaje) \\
\midrule
\textbf{Actores} & Jefe tarjador \este{(iniciador)} \\
\midrule
\textbf{Propósito} & Dejar constancia fotográfica y descriptiva de toda anomalía detectada en la
carga, de modo que el terminal pueda \textbf{acreditar el momento y las condiciones en que se
produjo} y se puedan determinar responsabilidades ante los seguros. \\
\midrule
\textbf{Tipo} & Primario $\cdot$ Esencial \\
\midrule
\textbf{Referencias cruzadas} & Funciones R4.9, R5.1 a R5.4, R6.4 $\cdot$ Casos de uso CU-12
\este{(lo extiende)}, CU-15, CU-17 \\
\midrule
\textbf{Precondiciones} & La faena está en ejecución y existe un lote o línea de carga en proceso
de registro. \\
\midrule
\textbf{Postcondiciones} & La incidencia queda registrada con su tipo, cantidad afectada,
descripción y al menos una fotografía exclusiva, asociada al lote y al cliente correspondiente. El
reporte final la destacará en rojo. \\

\bottomrule

\end{longtable}
}

{\estilotabla
\begin{longtable}{C{0.04\textwidth} L{0.44\textwidth} L{0.44\textwidth}}

\caption{CU-13: curso normal de los eventos}
\label{tab:cu13_normal}
\\

\toprule
\textbf{\#} & \textbf{Acción de los actores} & \textbf{Respuesta del sistema} \\
\midrule
\endfirsthead

\toprule
\textbf{\#} & \textbf{Acción de los actores} & \textbf{Respuesta del sistema} \\
\midrule
\endhead

1 & El Jefe tarjador detecta una anomalía y separa físicamente la unidad afectada del resto del
lote. & \\
\midrule
2 & El Jefe tarjador selecciona la opción de registrar incidencia sobre el lote en curso. & \\
\midrule
3 & & El sistema presenta los tipos de incidencia disponibles: daño, faltante, sobrante y mal
embalaje. \\
\midrule
4 & El Jefe tarjador selecciona el tipo, indica la cantidad afectada y describe lo observado. & \\
\midrule
5 & & El sistema crea la incidencia, la asocia a la línea de carga y al cliente, y \textbf{exige al
menos una fotografía exclusiva}. \\
\midrule
6 & El Jefe tarjador toma una o más fotografías dedicadas de la unidad afectada. & \\
\midrule
7 & & El sistema comprime, sella y sube cada fotografía, asociándolas a la incidencia
\este{(CU-15)}. \\
\midrule
8 & & El sistema confirma el registro y devuelve al operario al flujo de separación
\este{(CU-12)}. \\

\bottomrule

\end{longtable}
}

{\estilotabla
\begin{longtable}{C{0.06\textwidth} L{0.38\textwidth} L{0.48\textwidth}}

\caption{CU-13: cursos alternos}
\label{tab:cu13_alt}
\\

\toprule
\textbf{Paso} & \textbf{Situación} & \textbf{Tratamiento} \\
\midrule
\endfirsthead

\toprule
\textbf{Paso} & \textbf{Situación} & \textbf{Tratamiento} \\
\midrule
\endhead

4 & La incidencia fue generada automáticamente por diferencia de cantidad (paso 5 de CU-12). &
El sistema precarga el tipo (\estado{FALTANTE} / \estado{SOBRANTE}) y la cantidad afectada; el
operario solo describe y fotografía. \\
\midrule

6 & El operario intenta continuar sin fotografiar. &
El sistema \textbf{bloquea el avance}: una incidencia sin evidencia fotográfica carece de valor
probatorio. \\
\midrule

7 & No hay conectividad. &
Las fotografías se persisten localmente y se encolan para sincronización \este{(CU-16)}, sin
interrumpir la faena. \\
\midrule

--- & El daño afecta a varios clientes del mismo contenedor. &
Se registra una incidencia por cada línea de carga afectada, cada una con su evidencia propia, de
modo que cada cliente reciba en su reporte únicamente lo que le concierne. \\

\bottomrule

\end{longtable}
}

\subsection{CU-18 --- Revisar y aprobar reporte}
\label{subsec:cu18}

{\estilotabla
\begin{longtable}{L{0.25\textwidth} L{0.69\textwidth}}

\caption{CU-18 Revisar y aprobar reporte: descripción expandida}
\label{tab:cu18_desc}
\\

\toprule
\textbf{Campo} & \textbf{Contenido} \\
\midrule
\endfirsthead

\toprule
\textbf{Campo} & \textbf{Contenido} \\
\midrule
\endhead

\textbf{Caso de uso} & CU-18 Revisar y aprobar reporte \\
\midrule
\textbf{Actores} & Supervisor de reportes \este{(iniciador)}, Jefe tarjador \este{(participante en
el curso alterno de recaptura)} \\
\midrule
\textbf{Propósito} & Controlar la calidad técnica y la pertinencia comercial de la evidencia antes
de que llegue al cliente final, evitando que se publiquen fotografías inutilizables o detalles
operativos internos que el terminal prefiere resolver antes de notificar. \\
\midrule
\textbf{Tipo} & Primario $\cdot$ Esencial \\
\midrule
\textbf{Referencias cruzadas} & Funciones R6.5 a R6.9 $\cdot$ Casos de uso CU-17, CU-19, CU-20,
CU-22 \\
\midrule
\textbf{Precondiciones} & Existe un reporte en estado \estado{BORRADOR}, generado automáticamente
al cierre de la faena (CU-17). El Supervisor está autenticado y posee el permiso de aprobación. \\
\midrule
\textbf{Postcondiciones} & El reporte queda en estado \estado{APROBADO}, con registro del usuario
que lo aprobó y la fecha y hora de aprobación, y habilitado para su envío al cliente (CU-20). \\

\bottomrule

\end{longtable}
}

{\estilotabla
\begin{longtable}{C{0.04\textwidth} L{0.44\textwidth} L{0.44\textwidth}}

\caption{CU-18: curso normal de los eventos}
\label{tab:cu18_normal}
\\

\toprule
\textbf{\#} & \textbf{Acción de los actores} & \textbf{Respuesta del sistema} \\
\midrule
\endfirsthead

\toprule
\textbf{\#} & \textbf{Acción de los actores} & \textbf{Respuesta del sistema} \\
\midrule
\endhead

1 & El Supervisor accede a la bandeja de reportes pendientes de revisión. & \\
\midrule
2 & & El sistema lista los reportes en estado \estado{BORRADOR} con su folio, faena, cliente(s) y
fecha de generación. \\
\midrule
3 & El Supervisor selecciona un reporte. & \\
\midrule
4 & & El sistema presenta el borrador del PDF: encabezado con fecha, hora y cuadrilla a cargo,
desglose de fotografías y detalle por cliente, con la mercancía dañada destacada en rojo. \\
\midrule
5 & El Supervisor revisa visualmente la totalidad de las fotografías y los datos. & \\
\midrule
6 & El Supervisor aprueba el reporte. & \\
\midrule
7 & & El sistema registra la aprobación (usuario, fecha y hora), cambia el estado del reporte a
\estado{APROBADO} y habilita su envío. \\
\midrule
8 & & El sistema notifica al Supervisor que el reporte está listo para ser despachado
\este{(CU-20)}. \\

\bottomrule

\end{longtable}
}

{\estilotabla
\begin{longtable}{C{0.06\textwidth} L{0.38\textwidth} L{0.48\textwidth}}

\caption{CU-18: cursos alternos}
\label{tab:cu18_alt}
\\

\toprule
\textbf{Paso} & \textbf{Situación} & \textbf{Tratamiento} \\
\midrule
\endfirsthead

\toprule
\textbf{Paso} & \textbf{Situación} & \textbf{Tratamiento} \\
\midrule
\endhead

5 & Una fotografía está borrosa, mal encuadrada o apuntando al cielo. &
El Supervisor la rechaza indicando el motivo \este{(CU-19)}. El sistema la marca como
\estado{RECHAZADA}, cambia el reporte a \estado{EN\_CORRECCION} y notifica al jefe tarjador para
que capture un reemplazo. Recibido este, el sistema regenera el PDF e incrementa su versión. \\
\midrule

5 & Una fotografía muestra un detalle menor que el terminal prefiere resolver internamente antes de
notificar al cliente. &
Mismo tratamiento anterior. Esta es precisamente la razón por la que se eliminó la visualización en
vivo por parte del cliente. \\
\midrule

5 & Falta evidencia de una etapa obligatoria. &
El sistema advierte la omisión e impide la aprobación hasta que la faena esté completa. \\
\midrule

5 & La faena presenta incidencias graves que ameritan revisión presencial. &
El Supervisor deja el reporte en \estado{EN\_CORRECCION} con una observación y escala el caso al
Jefe de faena. \este{(Flujo administrativo, parcialmente fuera del sistema.)} \\
\midrule

7 & El Supervisor no tiene permiso de aprobación. &
El sistema deniega la acción y registra el intento en la bitácora. \\

\bottomrule

\end{longtable}
}

\subsection{CU-20 --- Enviar reporte al cliente}
\label{subsec:cu20}

{\estilotabla
\begin{longtable}{L{0.25\textwidth} L{0.69\textwidth}}

\caption{CU-20 Enviar reporte al cliente: descripción expandida}
\label{tab:cu20_desc}
\\

\toprule
\textbf{Campo} & \textbf{Contenido} \\
\midrule
\endfirsthead

\toprule
\textbf{Campo} & \textbf{Contenido} \\
\midrule
\endhead

\textbf{Caso de uso} & CU-20 Enviar reporte al cliente \\
\midrule
\textbf{Actores} & Supervisor de reportes \este{(iniciador)}, Servidor de correo \este{(sistema
externo)}, Cliente del depósito \este{(receptor)} \\
\midrule
\textbf{Propósito} & Hacer llegar al importador, de forma inmediata y trazable, el expediente
completo y aprobado del estado en que se recibió su mercancía. \\
\midrule
\textbf{Tipo} & Primario $\cdot$ Esencial \\
\midrule
\textbf{Referencias cruzadas} & Funciones R6.10, R6.11 $\cdot$ Casos de uso CU-02, CU-18, CU-21 \\
\midrule
\textbf{Precondiciones} & El reporte está en estado \estado{APROBADO}. Los clientes involucrados
tienen al menos un correo destinatario registrado (CU-02). \\
\midrule
\textbf{Postcondiciones} & El reporte queda en estado \estado{ENVIADO}, con registro de fecha,
hora, destinatarios y acuse de envío. \\

\bottomrule

\end{longtable}
}

{\estilotabla
\begin{longtable}{C{0.04\textwidth} L{0.44\textwidth} L{0.44\textwidth}}

\caption{CU-20: curso normal de los eventos}
\label{tab:cu20_normal}
\\

\toprule
\textbf{\#} & \textbf{Acción de los actores} & \textbf{Respuesta del sistema} \\
\midrule
\endfirsthead

\toprule
\textbf{\#} & \textbf{Acción de los actores} & \textbf{Respuesta del sistema} \\
\midrule
\endhead

1 & El Supervisor pulsa el botón de envío sobre un reporte aprobado. & \\
\midrule
2 & & El sistema obtiene los destinatarios registrados de los clientes involucrados en la faena. \\
\midrule
3 & & El sistema compone el correo con su asunto y adjunta el PDF aprobado. \\
\midrule
4 & & El sistema entrega el mensaje al servidor de correo. \\
\midrule
5 & & El servidor de correo despacha el mensaje al Cliente del depósito. \\
\midrule
6 & & El sistema recibe el acuse, registra el envío (fecha, hora, destinatarios, estado
\estado{ENVIADO}) y cambia el estado del reporte. \\
\midrule
7 & & El sistema confirma el envío al Supervisor. \\

\bottomrule

\end{longtable}
}

{\estilotabla
\begin{longtable}{C{0.06\textwidth} L{0.38\textwidth} L{0.48\textwidth}}

\caption{CU-20: cursos alternos}
\label{tab:cu20_alt}
\\

\toprule
\textbf{Paso} & \textbf{Situación} & \textbf{Tratamiento} \\
\midrule
\endfirsthead

\toprule
\textbf{Paso} & \textbf{Situación} & \textbf{Tratamiento} \\
\midrule
\endhead

2 & Un cliente no tiene destinatarios de correo registrados. &
El sistema advierte la omisión y solicita completar la ficha del cliente (CU-02) antes de
enviar. \\
\midrule

4--5 & El servidor de correo no responde o rechaza el mensaje. &
El sistema registra el envío con estado \estado{FALLIDO} y el mensaje de error, programa un
reintento automático y notifica al Supervisor. \\
\midrule

--- & El PDF excede el tamaño máximo admitido por el servidor de correo. &
El sistema envía un enlace de descarga en lugar del adjunto. \este{(Comportamiento inferido;
requiere validación.)} \\
\midrule

--- & El cliente solicita el reporte tiempo después. &
Se recupera desde la consulta histórica (CU-21) sin necesidad de reenvío. \\

\bottomrule

\end{longtable}
}

\section{Casos de uso pendientes de expandir}
\label{sec:cu_pendientes}

\begin{pendiente}
Los siguientes casos de uso se documentaron únicamente en su forma narrativa de alto nivel
(Sección~\ref{sec:cu_narrativos}). Se dejan las plantillas listas para expandirlos si la pauta o el
profesor lo requieren. Se sugiere priorizar \textbf{CU-11} y \textbf{CU-14}, que completan el ciclo
de la faena en terreno, y \textbf{CU-01}, como representante del módulo de administración.
\end{pendiente}

{\estilotabla
\begin{longtable}{L{0.34\textwidth} C{0.16\textwidth} L{0.44\textwidth}}

\caption{Casos de uso pendientes de expansión y prioridad sugerida}
\label{tab:cu_pendientes}
\\

\toprule
\textbf{Caso de uso} & \textbf{Prioridad sugerida} & \textbf{Justificación} \\
\midrule
\endfirsthead

\toprule
\textbf{Caso de uso} & \textbf{Prioridad sugerida} & \textbf{Justificación} \\
\midrule
\endhead

CU-11 Registrar apertura de contenedor & Alta &
Completa el ciclo de terreno junto a CU-12 y CU-14; contiene la verificación del sello, de alto
valor probatorio. \\
\midrule

CU-14 Registrar cierre de faena & Alta &
Dispara la generación automática del reporte y valida la completitud de la evidencia. \\
\midrule

CU-01 Gestionar usuarios, roles y permisos & Media &
Representa el módulo de administración de maestros. \\
\midrule

CU-07 Registrar carga en modalidad tarja ciega & Media &
Modalidad de contingencia con flujo propio diferenciado. \\
\midrule

CU-21 Consultar faenas históricas y trazabilidad & Baja &
Caso de uso de consulta, de menor complejidad algorítmica. \\
\midrule

CU-06, CU-09, CU-10, CU-15, CU-16, CU-17, CU-19, CU-22 & Baja &
Casos de uso secundarios o ya descritos dentro del curso de eventos de los casos expandidos. \\

\bottomrule

\end{longtable}
}

\textbf{Plantilla para la expansión} \este{(copiar y completar)}.

{\estilotabla
\begin{longtable}{L{0.25\textwidth} L{0.69\textwidth}}

\caption{Plantilla en blanco para la expansión de un caso de uso}
\label{tab:cu_plantilla}
\\

\toprule
\textbf{Campo} & \textbf{Contenido} \\
\midrule
\endfirsthead

\toprule
\textbf{Campo} & \textbf{Contenido} \\
\midrule
\endhead

\textbf{Caso de uso} & \\
\midrule
\textbf{Actores} & \\
\midrule
\textbf{Propósito} & \\
\midrule
\textbf{Tipo} & \\
\midrule
\textbf{Referencias cruzadas} & \\
\midrule
\textbf{Precondiciones} & \\
\midrule
\textbf{Postcondiciones} & \\

\bottomrule

\end{longtable}
}

{\estilotabla
\begin{longtable}{C{0.04\textwidth} L{0.44\textwidth} L{0.44\textwidth}}

\caption{Plantilla en blanco para el curso normal de los eventos}
\label{tab:cu_plantilla_normal}
\\

\toprule
\textbf{\#} & \textbf{Acción de los actores} & \textbf{Respuesta del sistema} \\
\midrule
\endfirsthead

\toprule
\textbf{\#} & \textbf{Acción de los actores} & \textbf{Respuesta del sistema} \\
\midrule
\endhead

1 & & \\
\midrule
2 & & \\

\bottomrule

\end{longtable}
}

{\estilotabla
\begin{longtable}{C{0.06\textwidth} L{0.38\textwidth} L{0.48\textwidth}}

\caption{Plantilla en blanco para los cursos alternos}
\label{tab:cu_plantilla_alt}
\\

\toprule
\textbf{Paso} & \textbf{Situación} & \textbf{Tratamiento} \\
\midrule
\endfirsthead

\toprule
\textbf{Paso} & \textbf{Situación} & \textbf{Tratamiento} \\
\midrule
\endhead

& & \\

\bottomrule

\end{longtable}
}
